\section{研究计划进度安排及预期目标}

\subsection{进度安排}

结合本课题的研究内容、技术路线以及毕业设计整体时间要求，研究计划拟按照如下进度安排开展：

\begin{enumerate}
    \item \textbf{2026 年 2 月 24 日至 2026 年 3 月 2 日：前期调研与准备阶段。} 集中开展相关文献阅读、资料搜集与开源代码调研，重点梳理 3DGS 基础方法、鱼眼相机建模与标定、多视图几何恢复、雷达点云处理以及 3DGEER、3DGUT 等面向通用相机或畸变相机的相关工作；同时完成实验环境搭建、数据采集方案设计、多个场景下的鱼眼相机视频采集，以及相关开源代码的筛选与整理，为后续研究实现提供基础。

    \item \textbf{2026 年 3 月 3 日至 2026 年 3 月 16 日：鱼眼相机建模与内参标定阶段。} 基于 OpenCV fisheye 模型完成鱼眼相机成像建模，采集多组包含 Charuco 标定板的鱼眼图像，完成相机内参标定，获得较为稳定的焦距、主点和畸变系数等参数，并对标定误差和参数稳定性进行初步分析。

    \item \textbf{2026 年 3 月 17 日至 2026 年 3 月 30 日：鱼眼相机位姿恢复阶段。} 在内参已知的条件下，研究更加稳定的鱼眼相机位姿恢复流程，完成特征提取、特征匹配、运动恢复结构与光束法平差优化，实现鱼眼图像的相机外参与稀疏点云恢复，并分析鱼眼模型参数过拟合对相机内外参估计精度和后续重建质量的影响。

    \item \textbf{2026 年 3 月 31 日至 2026 年 4 月 13 日：鱼眼图拆分与 3DGS 基线构建阶段。} 基于已标定的鱼眼相机参数，研究面向射线空间采样的鱼眼图拆分方法，完成从单张鱼眼图像到多张虚拟透视图的重采样流程，并由鱼眼相机外参与虚拟透视相机相对姿态直接计算透视图外参；在此基础上，将拆分得到的透视图接入传统 3DGS 训练管线，完成面向鱼眼输入的 3DGS 基线系统搭建。

    \item \textbf{2026 年 4 月 14 日至 2026 年 4 月 27 日：透视图外参优化与高斯增密改进阶段。} 在完成鱼眼图拆分后，进一步对透视图执行特征提取、特征匹配与光束法平差优化，以继续精化透视图外参；同时研究更优的高斯增密策略，提升场景细节区域的学习能力，并初步验证鱼眼图拆分策略、外参精化流程以及高斯增密改进对重建质量的影响。

    \item \textbf{2026 年 4 月 28 日至 2026 年 5 月 4 日：雷达点云筛选与辅助初始化阶段。} 完成原始雷达点云的筛选、噪声抑制与离群点剔除，研究利用雷达点云进行高斯位置初始化的方法，并探索利用逐帧雷达数据为鱼眼相机外参恢复提供更稳定初值的实现流程，为后续雷达几何约束引入奠定基础。

    \item \textbf{2026 年 5 月 5 日至 2026 年 5 月 11 日：雷达几何约束构建与系统迭代完善阶段。} 基于雷达点云尝试恢复场景粗表面表示，并将其转化为 3DGS 优化中的几何约束；完成鱼眼图拆分、透视图外参优化、高斯增密与雷达约束的联合系统集成，并在 2026 年 5 月 11 日前完成主要算法与系统功能的迭代改进，为后续论文定稿预留充足时间。

    \item \textbf{2026 年 5 月 12 日至 2026 年 5 月 18 日：实验评估与毕业论文完成阶段。} 集中开展系统实验与结果整理，完成与 3DGEER、3DGUT 等方法的量化对比，并采用 PSNR、SSIM 和 LPIPS 三项指标进行评估；同时完成消融实验，分析鱼眼图拆分方式、透视图外参进一步优化流程、高斯增密策略以及雷达约束设计对结果的影响。在此基础上完成毕业论文撰写、修改与定稿，确保于 2026 年 5 月 18 日前完成毕业论文。

    \item \textbf{2026 年 5 月 19 日至 2026 年 6 月 2 日：答辩准备与完善阶段。} 根据论文提交后的反馈继续完善文字表述、实验材料和图表展示，完成毕业答辩演示文稿制作、答辩讲稿整理以及多轮答辩演练，确保于 2026 年 6 月 2 日顺利完成毕业答辩。
\end{enumerate}

\subsection{预期目标}

结合本课题的研究内容与技术路线，预期达到以下目标：

\begin{enumerate}
    \item 完成基于 OpenCV fisheye 模型的鱼眼相机成像建模与标定流程，实现利用 Charuco 标定板进行鱼眼相机内参标定，并建立较为稳定的特征提取、匹配与光束法平差位姿恢复流程，获得较准确的相机内外参。

    \item 完成面向鱼眼输入的 3DGS 场景重建流程设计与实现。具体包括基于射线空间采样的鱼眼图拆分方法、由鱼眼相机外参推导透视图外参的方法，以及对拆分后透视图进一步进行特征提取、匹配和光束法平差优化的流程，使系统能够兼容传统 3DGS 训练管线并获得较高质量的场景重建结果。

    \item 完成雷达点云筛选与几何约束构建方法研究，探索将雷达信息用于高斯位置初始化、相机外参初始化以及场景表面几何约束构建的可行方案，从而提升低纹理区域和视角覆盖不足条件下的重建稳定性与跨视角几何一致性。

    \item 完成鱼眼 + 雷达室内场景重建系统的实验验证，建立与 3DGEER、3DGUT 等相关方法的对比实验，并采用 PSNR、SSIM 和 LPIPS 三项指标进行量化评价；同时通过消融实验分析鱼眼图拆分策略、透视图外参优化流程、高斯增密策略和雷达约束设计对最终结果的影响。

    \item 完成毕业论文撰写与答辩准备工作，形成较完整的研究实现、实验分析与结论总结，为顺利完成本科毕业设计与毕业答辩提供支撑。
\end{enumerate}

\subsection{已取得的成果}

截至目前，本课题已经完成了若干前期阶段性工作，并取得了以下成果：

\begin{enumerate}
    \item 已完成 Insta360 X5 鱼眼相机的标定工作，建立了基于 Charuco 标定板的鱼眼相机内参标定流程，并获得了可用于后续重建任务的相机内参结果。

    \item 已基于 COLMAP 完成相机外参标定，初步打通了从鱼眼图像输入到相机位姿恢复的处理流程，为后续鱼眼图拆分、透视图构建与重建实验提供了可靠的外参基础。

    \item 已完成鱼眼图像拆分为透视图的算法实现，初步打通了从鱼眼图像到多张虚拟透视图的重采样处理流程；在此基础上，已基于开源项目 LichtFeld-Studio 对 3DGS 训练进行了初步尝试，验证了拆分后数据与相机参数接入 3DGS 管线的可行性，并为后续训练策略调整与系统实现积累了实验经验。
\end{enumerate}
